\begin{thebibliography}{99}

\bibitem{DavenportSchmidt1967}
H.~Davenport and W.~M.~Schmidt,
\emph{Approximation to real numbers by quadratic \mbox{irrationals}},
Acta Arithmetica \textbf{13} (1967), 169--176.
\href{https://doi.org/10.4064/aa-13-2-169-176}{doi:10.4064/aa-13-2-169-176}.
Primary article consulted, including its minimal-point and determinant argument.

\bibitem{Lehmer1964}
D.~H.~Lehmer, \emph{On a problem of St\o rmer},
Illinois Journal of Mathematics \textbf{8} (1964), 57--79.
\href{https://doi.org/10.1215/ijm/1256067456}{doi:10.1215/ijm/1256067456}.
The introduction describes St\o rmer's original coefficient method and
smooth-coordinate lemma.  The sharper Lehmer index bound is not used here.

\bibitem{Stormer1897}
C.~St\o rmer, \emph{Quelques th\'eor\`emes sur l'\'equation de Pell
$x^2-Dy^2=\pm1$ et leurs applications},
Videnskabsselskabets Skrifter, I. Matematisk-Naturvidenskabelig Klasse,
Christiania, no.~2 (1897), 48 pp.  Historical original.

\bibitem{Klazar2010}
M.~Klazar, \emph{St\o rmer's solution of the unit equation $x-y=1$},
author's exposition, 25 August 2010.
\url{https://kam.mff.cuni.cz/~klazar/stormer.pdf}.
Consulted as an exposition of the elementary prime-index proof.

\bibitem{Luckhardt1989}
H.~Luckhardt, \emph{Herbrand-Analysen zweier Beweise des Satzes von Roth:
Polynomiale Anzahlschranken},
Journal of Symbolic Logic \textbf{54} (1989), 234--263.
\href{https://doi.org/10.2307/2275028}{doi:10.2307/2275028}.

\bibitem{KohlenbachHerbrand}
U.~Kohlenbach, \emph{Herbrand's theorem and extractive proof theory},
author's manuscript, 1 September 2008, especially Theorem~1.9.
\url{https://www2.mathematik.tu-darmstadt.de/~kohlenbach/Gazette.pdf}.

\bibitem{AdiceamShirandami2026}
F.~Adiceam and V.~Shirandami,
\emph{Probabilistic effectivity in the Subspace Theorem},
Research in Number Theory \textbf{12}, article 8 (2026).
Published 13 January 2026.
\href{https://doi.org/10.1007/s40993-025-00692-0}{doi:10.1007/s40993-025-00692-0}.
The introduction describes the outstanding effective-height problem;
its probabilistic results are not a uniform cutoff for individual inputs.

\bibitem{ProveItPR}
V.~Reshetnikov and contributors, ProveIt,
\texttt{ConciseIPRA/PR.lean}, inspected at commit
\texttt{b33ee1acb02b7619dfef419e61c3e343735a7bc4},
15 November 2026.
\href{https://github.com/VladimirReshetnikov/ProveIt/blob/b33ee1acb02b7619dfef419e61c3e343735a7bc4/Logic/PrimitiveRecursiveArithmetic/ConciseIPRA/Lean/ConciseIPRA/PR.lean}{Pinned source: PR syntax, terms, axioms, and quantifier-free induction}.

\bibitem{ProveItDeduction}
V.~Reshetnikov and contributors, ProveIt,
\texttt{ConciseIPRA/Deduction.lean}, same commit.
\href{https://github.com/VladimirReshetnikov/ProveIt/blob/b33ee1acb02b7619dfef419e61c3e343735a7bc4/Logic/PrimitiveRecursiveArithmetic/ConciseIPRA/Lean/ConciseIPRA/Deduction.lean}{Pinned source: intuitionistic natural deduction and equality}.

\bibitem{ProveItBi}
V.~Reshetnikov and contributors, ProveIt,
\texttt{ConciseIPRA/BiInterpretation.lean}, same commit.
\href{https://github.com/VladimirReshetnikov/ProveIt/blob/b33ee1acb02b7619dfef419e61c3e343735a7bc4/Logic/PrimitiveRecursiveArithmetic/ConciseIPRA/Lean/ConciseIPRA/BiInterpretation.lean}{Pinned source: bi-interpretation and theorem transport}.

\end{thebibliography}
